\medskip

Un jeu vidéo récompense par un objet tiré au sort les joueurs ayant remporté un défi. L'objet tiré peut être \og commun \fg{} ou \og rare \fg. Deux types d'objets communs ou rares sont disponibles, des épées et des boucliers.

Les concepteurs du jeu vidéo ont prévu que :

\begin{itemize}
\item la probabilité de tirer un objet rare est de 7\,\% ;
\item si on tire un objet rare, la probabilité que ce soit une épée est de 80\,\% ;
\item si on tire un objet commun, la probabilité que ce soit une épée est de 40\,\%.
\end{itemize}

\emph{Les parties \textbf{\rm A} et \textbf{\rm B} sont indépendantes.}

\bigskip

\textbf{Partie A}

\medskip

Un joueur vient de remporter un défi et tire au sort un objet. On note :

\begin{itemize}[label=$\bullet~$]
\item $R$ l'évènement \og le joueur tire un objet rare\fg{} ;
\item $E$ l'évènement \og le joueur tire une épée \fg{} ;
\item $\overline{R}$ et $\overline{E}$ les évènements contraires des évènements $R$ et $E$.
\end{itemize}

\medskip

\begin{enumerate}
\item Dresser un arbre pondéré modélisant la situation, puis calculer $P(R \cap E)$.
\item Calculer la probabilité de tirer une épée.
\item Le joueur a tiré une épée. Déterminer la probabilité que ce soit un objet rare. Arrondir
le résultat au millième.
\end{enumerate}

\bigskip

\textbf{Partie B}

\medskip

Un joueur remporte $30$ défis.

On note $X$ la variable aléatoire correspondant au nombre d'objets rares que le joueur obtient après avoir remporté $30$ défis. Les tirages successifs sont considérés comme indépendants.

\medskip

\begin{enumerate}
\item Déterminer, en justifiant, la loi de probabilité suivie par la variable aléatoire $X$. Préciser ses paramètres, ainsi que son espérance.
\item Déterminer $P(X < 6)$. Arrondir le résultat au millième.
\item Déterminer la plus grande valeur de $k$ telle que $P(X \geqslant k) \geqslant 0,5$. Interpréter le résultat dans le contexte de l'exercice.
\item Les développeurs du jeu vidéo veulent proposer aux joueurs d'acheter un \og ticket
d'or \fg{} qui permet de tirer $N$ objets. La probabilité de tirer un objet rare reste de 7\,\%.

Les développeurs aimeraient qu'en achetant un ticket d'or, la probabilité qu'un joueur obtienne au moins un objet rare lors de ces $N$ tirages soit supérieure ou égale à $0,95$.

Déterminer le nombre minimum d'objets à tirer pour atteindre cet objectif. On veillera à détailler la démarche mise en œuvre.
\end{enumerate}
